\part{Lemma}

\chapter{Algebra}

\begin{lemma}
  \label{lem:interpolation}
  For any polynomial \( A \in \CC[X] \) of degree \( n-1 \), any point set \( \{p_0, \ldots, p_{n-1}\} \) of size \( n \), \( \{(p_i, A(p_i))\}_{i=0}^{n-1} \) determines \( A(x) \) and vice versa.
\end{lemma}

\chapter{Analysis}

\begin{lemma}[\textbf{Master Theorem (simplified)}]
  \label{lem:master-sim}
  Let the recurrence relation be:
  \[ T(n) = aT\left(\frac{n}{b}\right) + n^d \]
  where $a \ge 1$ is the number of subproblems, $b > 1$ is the factor by which the problem size is reduced, and $d \ge 0$. The asymptotic bound of $T(n)$ is given by:

  \begin{enumerate}
      \item If $\log_b a > d$, then $T(n) = O(n^{\log_b a})$.
      \item If $\log_b a = d$, then $T(n) = O(n^d \log n)$.
      \item If $\log_b a < d$, then $T(n) = O(n^d)$.
  \end{enumerate}
\end{lemma}

\begin{lemma}[\textbf{Master Theorem}]
  \label{lem:master}
  Let $a \ge 1$ and $b > 1$ be constants, let $f(n)$ be a function, and let $T(n)$ be defined on the nonnegative integers by the recurrence:
  \[ T(n) = aT(n/b) + f(n) \]

  Then $T(n)$ has the following asymptotic bounds:
  \begin{itemize}
      \item \textbf{Case 1:} If $f(n) = O(n^{\log_b a - \epsilon})$ for some constant $\epsilon > 0$, \\
      then $T(n) = \Theta(n^{\log_b a})$.
      
      \item \textbf{Case 2:} If $f(n) = \Theta(n^{\log_b a} \log^k n)$ for some constant $k \ge 0$, \\
      then $T(n) = \Theta(n^{\log_b a} \log^{k+1} n)$.
      
      \item \textbf{Case 3:} If $f(n) = \Omega(n^{\log_b a + \epsilon})$ for some constant $\epsilon > 0$, \\
      and if $a f(n/b) \le c f(n)$ for some constant $c < 1$ and all sufficiently large $n$ (Regularity Condition), \\
      then $T(n) = \Theta(f(n))$.
  \end{itemize}
\end{lemma}
